% Tiril ar minya rimbë
\title{FAIRa Telemnar umë Cenda-Data Mindonar}
\subtitle{Ranta 1: Lo Carolo Taurë na ing.grid}
\shorttitle{FAIRa Telemnar}

% Lemba 1
\FDISetText{1}{1}{1}{Lo Carolo Taurë tenna Carolo Pitya.\\ Lo Karolingwa Minnascriptë ter Endor-ríki tenna \hspace*{0.75cm} NFDI4ING*. Cenda-data-mindon ná ar nauva sívë úvanima quenta.}
\FDISetText{1}{1}{2}{Númen-Francaríki}
\FDISetText{1}{1}{3}{Endor-ríki}
\FDISetText{1}{1}{4}{Rómen-Francaríki}
\FDISetText{1}{1}{5}{Forostar}
\FDISetText{1}{1}{6}{Baltëa\\Eär}
\FDISetText{1}{1}{7}{Aran.\\Anglia}
\FDISetText{1}{1}{8}{Írlanda}
\FDISetText{1}{1}{9}{Atalantëa\\Eär}
\FDISetText{1}{1}{10}{Aran. Asturia}
\FDISetText{1}{1}{11}{Brussel}
\FDISetText{1}{1}{12}{Aachen}
\FDISetText{1}{1}{13}{Luxem\FDIStyledText{8pt}{2pt}{-}\\[-2pt]burg}
\FDISetText{1}{1}{14}{Verdun}
\FDISetText{1}{1}{15}{Darmstadt}
\FDISetText{1}{1}{16}{Strasburg}
\FDISetText{1}{1}{17}{Roma}
\FDISetText{1}{1}{18}{Emirat Cordova}
\FDISetText{1}{1}{19}{Ríki Fez}
\FDISetText{1}{1}{20}{Endorëa Eär}
\FDISetText{1}{1}{21}{Benevent}
\FDISetText{1}{1}{22}{Serbia}
\FDISetText{1}{1}{23}{Slavëa Lier}
\FDISetText{1}{1}{24}{* NFDI4ING ná i omentië ten ingeniar-vísindar mi Nationalë Cenda-Data-Mindon (NFDI)}
\FDISetText{1}{1}{25}{Limes Sorabicus}


% Lemba 2
\FDISetText{2}{1}{1}{Mi 8/9. loa-randa Carolo Taurë yestëa carëa voronda ista-mindonënen essëa minya scriptë-cuilë. Sandasë turëan nérëatan scriptorion mi nildorion márar ar mi aranion coamas mi Aachen.}
\FDISetText{2}{2}{1}{Exë lambë\\istien ná ve\\namba fëa\\harya.}
\FDISetText{2}{3}{1}{Carolo immo úpollë tecë. Írë tecëssë parmari, antanë er i quanta-anna i Karolus-monogrammanna.}
\FDISetText{2}{4}{1}{Nan Carolo tultanë istië dictiënen, enyatien ar panta-lehtien parmari sauronissen.}
\FDISetText{2}{4}{2}{Parmar et-\\istalë-lanca!}
\FDISetText{2}{5}{1}{An quetilë ar eressëa-harmë parmaron mi túra ríki, carnanë i \enquote{Karolingwa Minnascriptë} ontaina.}
\FDISetText{2}{5}{2}{Er pitya tengwar (minnascriptë)}
\FDISetText{2}{5}{3}{Senda, panta, quilda tengwa-canta}
\FDISetText{2}{5}{4}{I \enquote{i} né en tecaina ú puntëo}
\FDISetText{2}{5}{5}{Canta-linya-axan an arya quetilë}
\FDISetText{2}{5}{6}{carolo ná úhanda}
\FDISetText{2}{6}{1}{I \enquote{s} né tenna\\i 12. loa-randa\\tecaina ve \enquote{anda s}}
\FDISetText{2}{6}{2}{I tengwa\\\enquote{j} úmë ëa mi 8.\\ar 9. loa-randa\\en}
\FDISetText{2}{6}{3}{Ú-langa tengwar an arya quetilë}
\FDISetText{2}{6}{4}{\FDIUseText{2}{5}{6}}
\FDISetText{2}{6}{5}{ta ia\\lá nanwa}

% Lemba 3
\FDISetText{3}{1}{1}{Yondor ar indyor Carolo hostaner sina analoga cenda-mindon parmalambëo ar quilëo palantírëa.}
\FDISetText{3}{1}{2}{Nossë Carolo}
\FDISetText{3}{1}{3}{(Cilmë)}
\FDISetText{3}{1}{4}{Carolo Taurë}
\FDISetText{3}{1}{5}{Pippin\\[-2.5pt]i\\[-2.5pt]Raica}
\FDISetText{3}{1}{6}{Carolo\\[-2.5pt]i\\[-2.5pt]Vinya}
\FDISetText{3}{1}{7}{Karlmann}
\FDISetText{3}{1}{8}{Ludwig\\[-2.5pt]i\\[-2.5pt]Aire}
\FDISetText{3}{1}{9}{Lothar I.}
\FDISetText{3}{1}{10}{Pippin}
\FDISetText{3}{1}{11}{Ludwig II.\\[-2.5pt]i\\[-2.5pt]Germanë}
\FDISetText{3}{1}{12}{Carolo II.\\[-2.5pt]i\\[-2.5pt]Pilha}
\FDISetText{3}{2}{1}{Minë indyo, \FDIUseText{3}{1}{9}, né mi\\loa 843 túra rimbëo mi Sáttë\\Verdunwa.}
\FDISetText{3}{3}{1}{Sina sáttë ná nórima ve nostalë i yestaron Francia ar Germania (\FDIUseText{1}{1}{2} ar \FDIUseText{1}{1}{4}).}
\FDISetText{3}{3}{2}{\FDIUseText{1}{1}{2}}
\FDISetText{3}{3}{3}{\FDIUseText{1}{1}{4}}
\FDISetText{3}{4}{1}{Lothar immo\\hirnë i\\Endor-ríki.\\Ta lendë\\ve anda,\\nalda\\rimba lo i\\Forostar tenna i\\Endorëa Eär.}
\FDISetText{3}{4}{2}{\FDIUseText{1}{1}{3}}
\FDISetText{3}{5}{1}{Sina \FDIUseText{1}{1}{3} carnëa i nostalë-erma\\Európëa omentiëo.}
\FDISetText{3}{6}{1}{Rimbë carmar i\\Európëa Omentiëo sí harir\\apa sina vanwa celma\\(Brussel, Luxemburg, Strasburg, Roma).}

% Lemba 4
\FDISetText{4}{1}{1}{Sívë martë i Comiciade\\mi Tecnología-mindon Europaplatz mi\\Aachen panta andor telemnar-meldor\\ilya ambaro.}
\FDISetText{4}{1}{2}{Harir ilyë\\quetë-huinar\\i vanima\\nómessë?}
\FDISetText{4}{1}{3}{Sina tië\\na DEUS\\AIX MACHINA-\\cesië!}
\FDISetText{4}{2}{1}{Ar mi C.A.R.L. RWTH Aachen\\carir entuluvar sinan yá.}
\FDISetText{4}{2}{2}{I rista\\lelya ar lelya\\na data-\\quenta!}
\FDISetText{4}{2}{3}{Ná\\Carolo Pitya\\en tulna?}
\FDISetText{4}{3}{1}{Mmmmhh....\\Lá\\raica...}
\FDISetText{4}{4}{1}{Toc}
\FDISetText{4}{5}{1}{Aiya\\Évariste!}
\FDISetText{4}{6}{1}{Á,\\sí nalyë,\\Alfred!}

% Lemba 5
\FDISetText{5}{1}{1}{Telemnar lya umë\\lúmë lyëo mi Cadil~Methodikë\\i\\Artificialë Intelligentiëo nányë melnë. Manen cenilyë, sí carëa i\\minya FAIRa telemnar ambardo?}
\FDISetText{5}{2}{1}{I minya\\FAIRa telemnar\\ambardo?}
\FDISetText{5}{3}{1}{Quetin FAIR mi sanwë\\cenda-data-mindono (FDI).}
\FDISetText{5}{3}{2}{I FAIR-axani\\nar sí:\\Findable, Accessible,\\Interoperable ar\\Re-usable.*}
\FDISetText{5}{3}{3}{*Hirima, minna-tulima, omentima ar ata-notima}
\FDISetText{5}{4}{1}{Ta mauya lyen\\nyarë nin\\sakon,\\Évariste.}
\FDISetText{5}{5}{1}{Aqua laicë. Lelyë acárië ISBN-nótië\\ambë i telemnar lyava,\\an sa ná laicë hirien\\mi parma-mart.}
\FDISetText{5}{6}{1}{Ná, lo ISBN-nótië má cenda\\i náma, lambë-nórë,\\quettarindo-nótië ar\\essë-nótië.}
\FDISetText{5}{6}{2}{Náma\\(parma)}
\FDISetText{5}{6}{3}{Quettarindo-nótië\\(Granus Verlag)}
\FDISetText{5}{6}{4}{Essë-\\nótië}
\FDISetText{5}{6}{5}{978-3-9822083-4-3  }
\FDISetText{5}{6}{6}{Nórë-nótië (Germania,\\Austria ar Helvetica)}
\FDISetText{5}{6}{7}{Tanca-\\nótië}

% Lemba 6
\FDISetText{6}{1}{1}{Maura ná ISBN-nótië, an minna\\i Lista Parmaron Hirimon (VLB)\\naitëa. Írë uas tanna, ná tanca urwa\\parma lya cambannar antien.}
\FDISetText{6}{2}{1}{I VLB ná i\\endëa, mirima\\data-sarna ar cenda-tólië\\an Germanë-lambëa\\parma-mart.}
\FDISetText{6}{3}{1}{Yára. Më istar\\notilmë DOI (Digital Object Identifier)\\enga ISBN-nótië.}
\FDISetText{6}{4}{1}{Ta ná eressëa,\\voronda essë-tanna\\náma ilqua,\\sívë digitalë ar\\hravanwa.}
\FDISetText{6}{4}{2}{DOI-lehtar}
\FDISetText{6}{4}{3}{Armar-sarnë}
\FDISetText{6}{4}{4}{https://doi.org/10.5281/zenodo.19468332}
\FDISetText{6}{4}{5}{Epessë}
\FDISetText{6}{4}{6}{Náma}
\FDISetText{6}{4}{7}{Hya ve QR-tengwë}
\FDISetText{6}{5}{1}{Minë meldo\\\hspace*{0.75cm} NFDI4INGo, i DataCite-omentië,\\anta DOI-r.*}
\FDISetText{6}{5}{2}{*DataCite ná register-omentië DOI-in}
\FDISetText{6}{6}{1}{DOI-r nar mi\\istalëa\\carmar i axan\\an tanca\\ettulë-tannë.}
\FDISetText{6}{6}{2}{Ar manen\\aquapa\\carir\\te?}

% Lemba 7
\FDISetText{7}{1}{1}{Te tanner\\digitalë námar.}
\FDISetText{7}{1}{2}{Digitalë náma 1}
\FDISetText{7}{1}{3}{Digitalë náma 2}
\FDISetText{7}{1}{4}{Digitalë náma 3}
\FDISetText{7}{1}{5}{Digitalë náma 4}
\FDISetText{7}{1}{6}{Digitalë náma 5}
\FDISetText{7}{1}{7}{Digitalë náma 6}
\FDISetText{7}{2}{1}{Maer. Sívë istalëa\\carmar voronwëa citima nar, ú\\liantëar \enquote{firëa} lelya, lá?}
\FDISetText{7}{3}{1}{Aqua sívë ná! Ilya digitalë náma,\\ná i quenta, data hya lírë-elen, harya\\ú-lumbëa tengwë.}
\FDISetText{7}{4}{1}{An i DOI ú\\vistaima ná, ta carë\\voronda tanna.}
\FDISetText{7}{5}{1}{Hmmm ...\\Manen má\\ta arya\\cénima carë?}
\FDISetText{7}{6}{1}{Aqua laicë.\\Më istar\\melmë notië\\i canwa-\\lótë-nyarna.}

% Lemba 8
\FDISetText{8}{1}{1}{Baunir}
\FDISetText{8}{1}{2}{Sívë\\or canwa-\\lótë\\i mirima lótëo\\nyarna ná,\\sívë mi Digitalë\\Náma (DO) i\\data-mirima\\nyarna ná.}
\FDISetText{8}{2}{1}{Sina canta tanna mai\\manen DO carna ná.}
\FDISetText{8}{2}{2}{Bit-raina}
\FDISetText{8}{2}{3}{Meta-data\\ \& \\context}
\FDISetText{8}{2}{4}{Operations}
\FDISetText{8}{2}{5}{Persistent Identifier}
\FDISetText{8}{3}{1}{Mi endë ná i bit-raina.\\Ta carna ná dataissen mi canta\\nolo ar minëo.}
\FDISetText{8}{4}{1}{Mi venya canwa-lótë-nyarna\\ta ná i mirima\\lótëo, i baunir.}
\FDISetText{8}{5}{1}{Apa sa nar i metadata.\\Te nyarir i mirima quettassen.\\Mi venya lúmë tanna ná\\ \enquote{\FDIUseText{8}{1}{1}}.}
\FDISetText{8}{5}{2}{\FDIUseText{8}{1}{1}}

% Lemba 9
\FDISetText{9}{1}{1}{I enta rindë ar operations\\nyarëa mana má carë\\i mirimainen.}
\FDISetText{9}{1}{2}{Operations}
\FDISetText{9}{2}{1}{NFDI4ING tirë\\an antëa ambë lá er download-\\carmë mi operations.}
\FDISetText{9}{3}{1}{I eterindë ná i Persistent\\Identifier (PID). PID ná i or-essë\\rimbë axanion, i carir digitalë\\námar hirima.}
\FDISetText{9}{3}{2}{PID}
\FDISetText{9}{3}{3}{Mi venya lúmë ta ná DOI.}
\FDISetText{9}{4}{1}{Nánë nar ar exë axanir, ve\\z.B. URN, Handle hya ORCID.*}
\FDISetText{9}{4}{2}{*URN = Uniform Resource Name Handle = Handle System ORCID = ten Open Researcher and Contributor ID}
\FDISetText{9}{5}{1}{Ilya DOI\\ná san PID\\nan lá ilya PID\\ná imya\\DOI.}
\FDISetText{9}{5}{2}{Ai!\\Tulyal nin\\sundo mir\\i canwa-\\lótë,\\Évariste!}
\FDISetText{9}{6}{1}{HA HA HA!\\Ná, tecin lyë\\alfabetëa sí,\\Alfred!}
\FDISetText{9}{6}{2}{Tecilyë\\nin\\alfabetëa?}

% Lemba 10
\FDISetText{10}{1}{1}{Mi FDI i quetta\\ \enquote{Alphabetisieren} nyarëa i lelya mi haryalë\\téra curur ar istië,\\i nar mauya, an cenda-data\\lesë ar tecë polë.}
\FDISetText{10}{2}{1}{Írë úquen polë te lesë,\\i arya armar-sarnë úanta ar\\i data quildë lemir.}
\FDISetText{10}{2}{2}{Ta ista Carolo Taurë yá.\\Sintes i mindon\\ú istiëo ná úanwa.}
\FDISetText{10}{3}{1}{I\\data-lambë\\mauya lesë\\parma.}
\FDISetText{10}{4}{1}{Carolo Taurë panyanë túra valta\\i aistari ar turuvar polir\\lesë ar tecë. Ar an\\te enyatir tancavë. Raicar\\data-leryalë né yá\\urco.}
\FDISetText{10}{5}{1}{Sívë quetien,\\copy-paste-raicar mi\\9. loa-randa.}
\FDISetText{10}{5}{2}{HA HA HA!}
\FDISetText{10}{5}{3}{Ta nyarin\\maira\\interoperabilitië\\vanwiëo ar\\síriëo!}
\FDISetText{10}{6}{1}{Mana ná\\sinna\\lala?}

% Lemba 11
\FDISetText{11}{1}{1}{Aiya Tobias. Më nar sí carëa\\FAIRa telemnar, ar\\ta ná sí anwavë alassëa.}
\FDISetText{11}{2}{1}{Ta ná tulca. Yestaurë antanen\\vinya paper nin ing.grid.\\Lelyë polir carë ta sina telemnaren\\yando.}
\FDISetText{11}{2}{2}{ing.grid?\\Mana ná\\ta?}
\FDISetText{11}{3}{1}{ing.grid ná\\diamond Open Access-\\journal NFDI4INGo.}
\FDISetText{11}{3}{2}{Tassë\\ëa yando\\preprint-\\sérvë.}
\FDISetText{11}{4}{1}{Ta ná alata sanwë, Tobias.\\I telemnar san nauva yá\\cénima or i preprint-sérvë.}
\FDISetText{11}{5}{1}{Aqua! Omentië-review nauva\\tassë i minya tyelca. Ta ná, i\\telemnar polë san quétaina\\ar metya.}
\FDISetText{11}{6}{1}{Maira!\\Masse panyan\\quanta-anna\\ninya?}

% Lemba 12
\FDISetText{12}{1}{1}{DARO! DARO! DARO! Minya mauya men\\i telemnar target-lambenna tulcë.\\Sina an Carolo Pitya yá\\harya tana-parma.}
\FDISetText{12}{2}{1}{Sina tana-parma\\hirima ná nu sina HANDLE.}
\FDISetText{12}{2}{2}{http://hdl.handle.net/2268.2/12918}
\FDISetText{12}{2}{3}{Ve citië, ta cénima ná sívë: \cite{SIMO21}}
\FDISetText{12}{3}{1}{Ma.\\San hlapilmë\\ta ve PDF\\am orta.}
\FDISetText{12}{4}{1}{Aaaaar...\\Á tulya!}
\FDISetText{12}{4}{2}{KLIK}
\FDISetText{12}{5}{1}{Parma lya ná alassëa ortaina.}
\FDISetText{12}{5}{2}{PLING!}
\FDISetText{12}{6}{1}{Sí nauva\\alassëa!}
\FDISetText{12}{6}{2}{Tyelë i\\minya ranto}

% Hannalë
\FDISetText{13}{1}{1}{The authors would like to thank the Federal Government and the Heads of Government of the Länder, as well as the Joint Science Conference (GWK), for their funding and support within the framework of the NFDI4ING consortium. Funded by the German Research Foundation (DFG) - project number 442146713 \cite{DFG22}.}